% Part: first-order-logic
% Chapter: models-theories
% Section: theories

\documentclass[../../../include/open-logic-section]{subfiles}

\begin{document}

\olfileid{fol}{mat}{the}

\olsection{د لومړۍ درجې تيوريو بېلګې}

\begin{ex}
په ژبه~$\Lang L_<$ کښې د سختو خطي ترتيبونو تيوري د لاندې سټ له خوا
په بديهي اصولو ټاکل کېږي:
\begin{align*}
\{\quad & \lforall[x][\lnot x < x], \\
& \lforall[x][\lforall[y][((x < y \lor y <
    x) \lor x = y)]], \\
& \lforall[x][\lforall[y][\lforall[z][((x < y
      \land y < z) \lif x < z)]]] \quad \}
\end{align*}
دا مطلوب !!{structure}s په بشپړ ډول رانغاړي: هر سخت خطي ترتيب د دې
بديهي نظام مدل دے؛ او برعکس، که~$R$ پر يوه سټ~$X$ خطي ترتيب وي، نو
هغه جوړښت~$\Struct M$ چې $\Domain M = X$ او
$\Assign{<}{M} = R$ لري، د دې تيورۍ مدل دے۔
\end{ex}

\begin{ex}
د ګروپونو تيوري په هغه ژبه کښې چې~$\Obj 1$ (!!{constant}) او~$\cdot$
(دوه‌ځاييزه !!{function}) لري، د لاندې فورمولونو له خوا په بديهي
اصولو ټاکل کېږي:
\begin{align*}
& \lforall[x][\eq[(x \cdot \Obj 1)][x]]\\
& \lforall[x][\lforall[y][\lforall[z][\eq[(x \cdot (y \cdot z))][((x
          \cdot y) \cdot z)]]]]\\
& \lforall[x][\lexists[y][\eq[(x \cdot y)][\Obj 1]]]
\end{align*}
\end{ex}

\begin{ex}
د پيانو د حساب تيوري د حساب په ژبه~$\Lang L_A$ کښې د لاندې جملو له
خوا په بديهي اصولو ټاکل کېږي۔
\begin{align*}
& \lforall[x][\lforall[y][(\eq[x'][y'] \lif \eq[x][y])]]\\
& \lforall[x][\eq/[\Obj 0][x']]\\
& \lforall[x][\eq[(x + \Obj 0)][x]]\\
& \lforall[x][\lforall[y][\eq[(x + y')][(x + y)']]]\\
& \lforall[x][\eq[(x \times \Obj 0)][\Obj 0]]\\
& \lforall[x][\lforall[y][\eq[(x \times y')][((x \times y) + x)]]]\\
& \lforall[x][\lforall[y][(x < y \liff \lexists[z][\eq[(z' + x)][y])]]]\\
\intertext{او ورسره د لاندې بڼې ټولې جملې}
& (!A(\Obj 0) \land \lforall[x][(!A(x) \lif !A(x'))]) \lif \lforall[x][!A(x)]
\end{align*}
ځکه چې د وروستۍ بڼې نامتناهي ډېرې جملې شته، دا بديهي نظام
نامتناهي دے۔ وروستۍ بڼه \emph{د استقرا شېما} بلل کېږي۔ (په اصل کښې
د استقرا شېما تر هغه لږه پېچلې ده چې دلته مو ښودلې ده۔)

وروستے بديهي اصل د~$<$ يو \emph{صريح تعريف} دے۔
\end{ex}

\begin{ex}
د خالصو سټونو تيوري د رياضۍ په بنسټونو (او فلسفه) کښې مهمه ونډه
لري۔ يو سټ هله خالص دے چې ټول !!{element}s ئې هم خالص سټونه وي۔ نو
تش سټ خالص ګڼل کېږي، خو هغه سټ چې کوم داسې څيز ئې !!a{element} وي
چې خپله سټ نۀ وي، خالص نۀ دے۔ په دې توګه خالص سټونه يوازې له تش سټ
څخه، بې له ``غيرسټي بنسټيزو عناصرو''، جوړ شوي سټونه دي؛ يعنې داسې
څيزونه نۀ لري چې خپله سټونه نۀ وي۔

لاندې فورمولونه د خالصو سټونو د يوې تيورۍ بديهي نظام ګڼل کېدے شي:
\begin{align*}
& \lexists[x][\lnot \lexists[y][y \in x]]\\
& \lforall[x][\lforall[y][(\lforall[z](z \in x \liff z \in y) \lif
\eq[x][y])]]\\
& \lforall[x][\lforall[y][\lexists[z][\lforall[u][(u \in z \liff
(\eq[u][x] \lor \eq[u][y]))]]]]\\
& \lforall[x][\lexists[y][\lforall[z][(z \in y \liff \lexists[u][(z \in
        u \land u \in x)])]]]\\
\intertext{او ورسره د لاندې بڼې ټولې جملې} &
\lexists[x][\lforall[y][(y \in x \liff !A(y))]]
\end{align*}
لومړے بديهي اصل وايي چې داسې يو سټ شته چې هېڅ !!{element}s نۀ لري
(يعنې~$\emptyset$ شته)؛ دويم وايي چې سټونه امتدادي دي؛ درېيم وايي چې د هر
دوو سټونو~$X$ او~$Y$ دپاره سټ~$\{X, Y\}$ شته؛ څلورم وايي چې د هر
سټ~$X$ دپاره سټ~$\cup X$ شته، چېرته چې~$\cup X$ د~$X$ د ټولو
عناصرو اتحاد دے۔

وروستۍ يادې شوې !!{sentence}s په ګډه \emph{د ساده ادراک شېما} بلل
کېږي۔ په اصل کښې وايي چې د هر~$!A(x)$ دپاره سټ
$\Setabs{x}{!A(x)}$ شته---نو په لومړي نظر يو رښتونے، ګټور او ښايي
حتا ضروري بديهي اصل ښکاري۔ ``ساده'' ځکه بلل کېږي چې په پايله کښې دا
تيوري د صدق وړ نۀ پرېږدي: که~$!A(y)$ د~$\lnot y \in y$ په توګه
واخلئ، لاندې !!{sentence} لاس ته راځي:
\[
\lexists[x][\lforall[y][(y \in x \liff \lnot y \in y)]]
\]
او دا !!{sentence} په هېڅ !!{structure} کښې نۀ صدق کوي۔
\end{ex}

\begin{ex}
د \emph{مېرېولوژۍ} يا جزپوهنې په ډګر کښې د \emph{برخه‌والي} اړيکه
بنسټيزه اړيکه ده۔ لکه د سټونو د تيوريو په شان، د برخه‌والي داسې
تيورۍ هم شته چې د دې اړيکې بېلابېل (او کله کله يو بل سره ټکر لرونکي)
تصورات په بديهي اصولو ټاکي۔

د مېرېولوژۍ ژبه يو دوه‌ځاييز اړوند سمبول~$\Obj P$ لري، او
$\Atom{\Obj P}{x, y}$ دا ``معنٰی'' لري چې~$x$ د~$y$ يوه برخه ده۔
کله چې دا تفسير په پام کښې ولرو، د دې ژبې !!a{structure} د
\emph{برخه‌والي جوړښت} بلل کېږي۔ البته د يوه دوه‌ځاييز اړوند هر
جوړښت په رښتيا د دې نوم وړ نۀ دے۔ د ``برخه‌والي'' د رانغاړلو دپاره
بايد~$\Assign{\Obj P}{M}$ ځينې شرطونه پوره کړي، او دا شرطونه د
برخه‌والي د تيورۍ د بديهي اصولو په توګه ټاکلے شو۔ مثلاً برخه‌والے
پر څيزونو جزوي ترتيب دے: هر څيز د خپل ځان (که څه هم \emph{غيرحقيقي})
برخه ده؛ دوه بېل څيزونه د يو بل برخې نۀ شي کېدے؛ او د يوه څيز د
برخې برخه خپله هم د هغه څيز برخه ده۔ پام وکړئ چې په دې معنٰی ``د
يوه څيز برخه ده'' له ``فرعي سټ دے'' سره ورته ده، خو له ``عنصر دے''
سره نۀ ده ورته، ځکه عنصرول نۀ انعکاسي دي او نۀ انتقالي۔
\begin{align*}
& \lforall[x][\Atom{\Obj P}{x,x}] \\
& \lforall[x][\lforall[y][((\Part{x}{y} \land \Part{y}{x})
      \lif \eq[x][y])]] \\
& \lforall[x][\lforall[y][\lforall[z][((\Part{x}{y} \land
        \Part{y}{z}) \lif \Part{x}{z})]]]\\
\intertext{سربېره پر دې، هر دوه څيزونه مېرېولوژيکي مجموعه لري (داسې
  څيز چې دغه دواړه څيزونه ئې برخې وي او په دې لحاظ تر ټولو کوچنے وي)۔}  &
\lforall[x][\lforall[y][\lexists[z][\lforall[u][(\Part{z}{u} \liff
        (\Part{x}{u} \land \Part{y}{u}))]]]]
\end{align*}
دا يوازې د برخه‌والي ځينې هغه بنسټيز اصول دي چې مابعدالطبيعتپوهان ئې
څېړي۔ خو نور اصول ډېر ژر داسې پېچلي شي چې تر ځينو تعريف شويو اړيکو
وړاندې ئې جوړول يا ليکل ستونزمن وي۔ مثلاً په مېرېولوژۍ کښې زياتره
دلچسپي لرونکي مابعدالطبيعتپوهان لاندې اصل هم معتبر ګڼي: هر کله چې
څيز~$x$ حقيقي برخه~$y$ ولري، نو داسې برخه~$z$ هم لري چې له~$y$ سره
هېڅ ګډه برخه ونۀ لري، او د~$y$ او~$z$ امتزاج~$x$ وي۔
\end{ex}

\end{document}
